% Submission build: MDPI's own class, Definitions/mdpi.cls, unmodified (see
% README.md for where Definitions/ came from). The stand-in article build is
% main_article.tex; both pull the same sections/ files.
%
% Build: "D:/environment/tools/tectonic/tectonic.exe" main.tex   (run from paper/)
%
% The template's first option is the journal; `mathematics' is in the class's
% journal list. MDPI's own template.tex passes no graphics driver, and the
% `pdftex' option that older MDPI templates carried must NOT be used here: it is
% handed down to article.cls, which then forces pdftex.def, and xcolor breaks on
% \pdfcolorstack under XeTeX. With no driver option, graphicx detects xetex.
\documentclass[mathematics,article,submit,moreauthors]{Definitions/mdpi}

% --- EPS logo under XeTeX -------------------------------------------------
% mdpi.cls includes Definitions/logo-mdpi.eps on the title page. XeTeX's
% xdvipdfmx cannot read EPS, so the include fails under tectonic. mdpi.cls is
% not edited; instead every .eps request is answered with the PDF of the same
% base name. Definitions/logo-mdpi.pdf is generated once from MDPI's own
% Definitions/logo-mdpi.eps -- see README.md for the command. Under a pdflatex
% build this block is harmless as long as the PDF is present; delete it if you
% build with pdflatex and prefer the EPS.
\makeatletter
\let\mdpi@includegraphics\includegraphics
\renewcommand{\includegraphics}[2][]{%
  \IfEndWith{#2}{.eps}%
    {\StrGobbleRight{#2}{4}[\mdpi@gfile]%
     \mdpi@includegraphics[#1]{\mdpi@gfile.pdf}}%
    {\mdpi@includegraphics[#1]{#2}}%
}
\makeatother
% --------------------------------------------------------------------------

% The class already loads amsmath, amssymb, amsthm, booktabs, graphicx, tikz,
% natbib, hyperref and cleveref, so none of those is loaded again here.
\usepackage{algorithm}
\usepackage{algpseudocode}
\usetikzlibrary{decorations.pathmorphing,positioning,arrows.meta}
\usepackage{pgfplots}
\pgfplotsset{compat=1.18}

% Theorem environments. mdpi.cls defines Theorem, Lemma, Corollary,
% Proposition, Definition, Remark, Example and proof in the house style; the
% section files were written against the lowercase amsthm names, so the two are
% tied together here rather than declared a second time. sections/06_theory.tex
% declares its own only when \theorem is still undefined, so this wins.
\makeatletter
\let\theorem\Theorem                 \let\endtheorem\endTheorem
\let\lemma\Lemma                     \let\endlemma\endLemma
\let\corollary\Corollary             \let\endcorollary\endCorollary
\let\proposition\Proposition         \let\endproposition\endProposition
\let\definition\Definition           \let\enddefinition\endDefinition
\let\remark\Remark                   \let\endremark\endRemark
\let\example\Example                 \let\endexample\endExample
\makeatother

% Marks a figure that comes from development data rather than from the sealed
% semesters. It expands to its argument; the submission version can redefine it.
\newcommand{\devnum}[1]{#1}
% Figures from the sealed terms, opened once after the protocol freeze.
\newcommand{\sealednum}[1]{#1}

% Pulls in a section if its file exists, and otherwise leaves a numbered
% placeholder so that section numbers and cross-references stay correct.
\newcommand{\papersection}[3]{%
  \IfFileExists{sections/#1.tex}%
    {\input{sections/#1}}%
    {\section{#2}\label{#3}%
     \textit{Not yet in the repository: the file
     \texttt{\detokenize{sections/#1.tex}} was not found.}}%
}

%=================================================================
% MDPI internal commands - do not modify
\firstpage{1}
\makeatletter
\setcounter{page}{\@firstpage}
\makeatother
\pubvolume{1}
\issuenum{1}
\articlenumber{0}
\pubyear{2026}
\copyrightyear{2026}
\datereceived{ }
\daterevised{ }
\dateaccepted{ }
\datepublished{ }

%=================================================================
% Front matter. Everything only the authors can supply is marked
% [to be completed by the authors]; no name, affiliation, e-mail or ORCID is
% invented here.
\Title{Prediction-Driven Non-Preemptive Scheduling with Bounded Overtaking for
Shared Execution Services under Deadline-Driven Bursty Load}
% This version of mdpi.cls (v6.5a, 2026-09-11) has no \TitleCitation and no
% \AuthorCitation; if the journal's current template adds them, the running-head
% forms go here.

% Two authors. Replace the bracketed fields; if both authors share one
% affiliation, delete address 2 and change the second author's mark to $^{1}$.
% ORCID (optional): add \orcidA{} / \orcidB{} after a name and define
% \newcommand{\orcidauthorA}{0000-0000-0000-0000} in the preamble.
\Author{[First Author Name] $^{1}$ and [Second Author Name] $^{2,}$*}
\AuthorNames{[First Author Name] and [Second Author Name]}

\address{%
$^{1}$ \quad [First author: department, institution, city and postcode, country]; [first author e-mail]\\
$^{2}$ \quad [Second author: department, institution, city and postcode, country]; [second author e-mail]}

\corres{Correspondence: [corresponding author e-mail]}

\abstract{A shared execution service runs jobs on a fixed pool of servers, without preemption, under a per-job time limit $L$. Reordering by predicted cost shortens most waits and can leave one job far behind. We prove a pathwise identity: for any non-preemptive work-conserving $k$-server policy and input, a job's delay relative to first-come first-served equals its net overtaken work divided by $k$, to within $(2-2/k)L$, an exact supremum. Bounding overtaken work is therefore sufficient for a per-job guarantee, and necessary up to an additive $O(kL)$; a count of overtakes is not necessary, and suffices only by charging $L$ per pass. A wrapper enforcing that budget bounds every wait against first-come first-served under arbitrary arrivals and predictions. It reads the budget only through a cap, so the shape below is free; the additive constant $(3-2/k)L$ is likewise exact, and the cap follows from a promise in seconds. Expected-cost ordering is an empirically motivated base policy rather than a multiserver optimality claim. We evaluate on two judging platforms, a serverless trace, two continuous-integration pools that check the simulator, and a compute farm. The empirical headline uses conservative outcome visibility: at the busiest development load, a \devnum{600}-second guard closes \devnum{0.429} of the FCFS-to-SJF tail-wait gap, with interval \devnum{[0.357, 0.533]}. Original-clock results are optimistic references. We report when the uniform certificate is too coarse for the observed waits and when cost concentration offers little benefit.}

% Source of the abstract's empirical values: outputs/dev_visibility/visibility_comparison.csv,
% level=2, Guard(600), gap_closed and gap_closed_lo/hi.
\keyword{non-preemptive scheduling; scheduling with predictions; multiserver queues; bounded overtaking; per-job guarantee; first-come first-served; runtime prediction}

% Mathematics prints MSC 2020 codes under the keywords: scheduling theory, performance
% evaluation and scheduling in computer systems, queueing theory, online algorithms.
\MSC{90B35; 68M20; 60K25; 68W27}

%=================================================================
\begin{document}

\maketitle

\papersection{01_introduction}{Introduction}{sec:intro}
\papersection{02_related_work}{Related Work}{sec:related}
\papersection{03_problem_model}{Problem Model}{sec:model}
\papersection{04_prediction}{Cost Prediction}{sec:prediction}
\papersection{05_scheduling}{Scheduling Algorithm}{sec:scheduling}
\papersection{06_theory}{Theory}{sec:theory}
\papersection{07_data}{Traces and Workload Characterisation}{sec:data}
\papersection{08_experiments}{Experiments}{sec:experiments}
% sections/09_limitations.tex ends with \section{Conclusions} (label sec:conclusions), so
% the Conclusions section MDPI's structure expects is typeset from that file and needs no
% \papersection call of its own.
\papersection{09_limitations}{Limitations and Threats to Validity}{sec:limits}

%=================================================================
% Back matter. Anything only the authors can supply is marked
% [to be completed by the authors] and is not invented.

% The supplementary file is paper/supplementary.tex, built on its own with the same class
% under the `supfile' article type and the same bibliography. Its S-items are listed here
% in the order they appear there.
\supplementary{The following supporting information can be downloaded at
\linksupplementary{s1}. Section~S1, cost prediction: the visibility audit of the feature
pipeline in full; Table~S1, predictors on the four non-judging traces by feature group;
the release delay treated as a sensitivity; the three predictor families and the
ablation ladder; Table~S2, predictors on the two judging traces; and the measured cost
of predicting. Section~S2, traces: the field audit of the main judging trace with the
release-statistics reconciliation and the two truncated archives; the
capacity-sharpness and circularity checks of the continuous-integration refit; what
supports the shared-pool design; and Table~S3, cost concentration with the provenance of
the analysis bound on every trace. Section~S3, experiments: Table~S4, the experimental
configuration and the selection protocol; Table~S5, the selected budget parameters at
the three promises with the nine per-shape grid winners; Table~S6, the age-relative and
queue-length ablation rows of Table~\ref{tab:guard}; Table~S7, the corrupted scores;
Table~S8, the residual of Theorem~\ref{thm:identity} on all nine cells; Table~S9, the
single-server configuration; Tables~S10--S12, the aging comparator and the package
predictor metrics; Table~S13, the alternative ranking losses; Table~S14, the
trace-construction sensitivities; and Figures~S1--S3. Section~S4, the deferred proofs:
Lemma~\ref{lem:gap} with its strictness, Theorem~\ref{thm:identity} with the two
directions of its extremes, and Proposition~\ref{prop:consistency} with the
two counterexamples to its converse, together with the remaining proofs and remarks
deferred from Section~\ref{sec:theory}. Section~S5, the reproducibility specification, including exact
grids, splits, information clocks and table provenance. Section~S6, the capacity sweep,
physical timed-work experiment, feedback counterexamples and LPC-EGEE counterfactual.
Section~S7, additional theory, with Table~S18 comparing the nearest guarantee
disciplines. Section~S8, the extremal constructions, each opening
with its own idea and then given in full with the tables of exact values it attains,
together with the achievable fraction at one server. Section~S9, the extended related
work: the per-paper detail behind Section~\ref{sec:related}, every work named there
being cited in the manuscript as well.}

\authorcontributions{[to be completed by the authors]. The template expects the
CRediT roles (conceptualization, methodology, software, validation, formal
analysis, investigation, data curation, writing---original draft preparation,
writing---review and editing, visualization, supervision, project
administration, funding acquisition), each attributed by author initials,
followed by ``All authors have read and agreed to the published version of the
manuscript.''}

\funding{[to be completed by the authors]. If there is no external funding, the
template asks for ``This research received no external funding''; otherwise the
grant number and the funding body go here.}

\institutionalreview{Not applicable.}

\informedconsent{Not applicable.}

\dataavailability{The public source datasets used here are available at the
locations cited, and we redistribute none of them. The CodeBench judging logs
are published by the Universidade Federal do Amazonas \cite{codebench_dataset};
the ACcoding dataset is published with its data descriptor \cite{accoding2024};
the Open University Learning Analytics dataset is published with its own
\cite{oulad2017}; the Azure Functions 2021 invocation trace is published in the
Azure Public Dataset repository \cite{azure2021trace}, which asks that
\cite{zhang2021harvested} be cited; and the Intel Netbatch logs are published in
the Parallel Workloads Archive \cite{netbatch2012trace}, which asks that
\cite{shai2014netbatch} be cited and that the logs be acknowledged to their
depositors. The LPC-EGEE log is also published in the Parallel Workloads Archive
\cite{lpc_egee_trace}, with acknowledgment requested for its data provider.
The continuous-integration trace of the two Firefox
hardware worker pools, was collected by us from Mozilla's public,
unauthenticated Taskcluster and Treeherder APIs
\cite{taskcluster_docs,treeherder_docs}; we do not redistribute it, and the
collection scripts that rebuild it from those APIs are part of the released
code. The Azure invocation trace is distributed under the publisher's Creative Commons
Attribution licence \cite{azure2021trace}; public availability does not imply an
absence of licence conditions. No restricted access agreement was used, and no
personal data were processed beyond the pseudonymous user identifiers the
published judging datasets already carry. The simulator, the predictors, the
guard implementation, configurations and archived development analyses are to be
released with the paper. Supplementary Section~S5 distinguishes the tables reproduced
by the package from the earlier development evidence. The release will be available at [repository URL to be
completed by the authors].}

\acknowledgments{[to be completed by the authors].}

\conflictsofinterest{[to be completed by the authors]. The template asks for
``The authors declare no conflicts of interest'' where that holds, and otherwise
for the competing interest and for a statement on the role of any funder in the
design, collection, analysis or interpretation, in the writing, or in the
decision to publish.}

\abbreviations{Abbreviations}{%
The following abbreviations are used in this manuscript:\\

\noindent
\begin{tabular}{@{}ll}
AUPRC & Area under the precision--recall curve\\
AUROC & Area under the receiver operating characteristic curve\\
CI & Continuous integration\\
FCFS & First-come first-served\\
GRU & Gated recurrent unit\\
MLP & Multilayer perceptron\\
OULAD & Open University Learning Analytics Dataset\\
RMSE & Root mean squared error\\
SJF & Shortest job first\\
SPJF & Shortest predicted job first\\
SRPT & Shortest remaining processing time
\end{tabular}
}

%=================================================================
% Ninth pass: the manuscript no longer carries an appendix. What stood in
% sections/A_proofs.tex was the idea of each extremal construction beside a pointer to
% the full argument; the ideas are now the opening of Supplementary Section~S8, beside
% those arguments, and Proposition A1 is Proposition 6 of Supplementary Section S8.5.
% sections/A_proofs.tex is kept as a stub recording where its content went, and is no
% longer \input, so \appendixtitles, \appendixstart and the \theH... anchor prefixes
% that went with it are gone as well.

%=================================================================
\begin{adjustwidth}{-\extralength}{0cm}

\reftitle{References}

% The class sets \bibliographystyle{Definitions/mdpi} itself; it is not set
% again here. The x-verified fields of refs.bib are read by no bst and do not
% print; only the short note fields (Accessed ..., Preprint arXiv:...) do.
\bibliography{../docs/related_work/refs}

\PublishersNote{}
\end{adjustwidth}
\end{document}
